\documentclass[11pt]{article}
\usepackage{amsmath,amsthm,amssymb}
\usepackage[margin=1.1in]{geometry}
\usepackage[numbers,sort&compress]{natbib}
\usepackage{hyperref}

\newtheorem{theorem}{Theorem}
\newtheorem{lemma}{Lemma}
\newtheorem{corollary}{Corollary}
\theoremstyle{remark}
\newtheorem{remark}{Remark}

\title{A sharp norm bound for weighted up-operators on the Boolean lattice}
\author{Gerald\thanks{Author name and affiliation to be supplied before submission.}
\and with assistance from an AI research agent\thanks{AI-assisted proof
verification and manuscript preparation; disclosure per journal policy.}}
\date{\today}

\begin{document}
\maketitle

\begin{abstract}
For $a=(a_1,\dots,a_m)\in\mathbb{R}^m$ and $1\le k\le m$, let
$U_k(a):\ell^2\binom{[m]}{k}\to\ell^2\binom{[m]}{k-1}$ be the weighted
up-operator (the weighted $(k,k-1)$ inclusion matrix)
\[
(U_k(a)f)(S)=\sum_{i\notin S} a_i\, f(S\cup\{i\}),\qquad |S|=k-1 .
\]
We prove the sharp inequality
$\|U_k(a)\|^2\le \frac{k(m-k+1)}{m}\sum_{i=1}^m a_i^2$,
with equality for all $a$ at the endpoints $k=1,m$ and, for interior $k$,
if and only if all $|a_i|$ are equal. The classical equal-weight identity
$\|U_k(\mathbf 1)\|^2=k(m-k+1)$ is the case $a=\mathbf 1$; the weighted
statement is not accessible to Johnson-scheme spectral methods, which fail
once the weights are generic. The proof is a coordinate-deletion induction
whose scalar closure defect factors as a perfect square.
\end{abstract}

\noindent\textbf{Keywords:} Boolean lattice, inclusion matrix, operator norm,
Johnson scheme, perfect-square factorization.\\
\noindent\textbf{MSC 2020:} 05E30, 15A60, 05A99.

\section{Introduction and statement of the result}

Let $m\ge 2$ and $1\le k\le m$. Write $[m]=\{1,\dots,m\}$, let
$\binom{[m]}{j}$ denote the $j$-subsets of $[m]$, and let
$\ell^2\binom{[m]}{j}$ be the Euclidean space with orthonormal basis
$\{e_S:|S|=j\}$. For weights $a=(a_1,\dots,a_m)\in\mathbb{R}^m$ define the
\emph{weighted up-operator}
$U_k(a):\ell^2\binom{[m]}{k}\to\ell^2\binom{[m]}{k-1}$ by
\[
(U_k(a)f)(S)=\sum_{i\notin S}a_i f(S\cup\{i\}),\qquad |S|=k-1 .
\]
In the subset basis $U_k(a)$ is the weighted $(k,k-1)$ inclusion matrix:
its $(S,T)$ entry is $a_i$ if $T=S\cup\{i\}$ and $0$ otherwise. It is the
weighted analogue of the ordinary inclusion (up) operator of the Boolean
lattice, whose spectral theory is classical
\cite{srinivasan2011,bannai1984,godsilroyle2001}; in particular
\begin{equation}\label{eq:classical}
\|U_k(\mathbf 1)\|^2=k(m-k+1).
\end{equation}
A natural question is whether coordinate weights can beat the
equal-weight spectral value: for which $C$ does
$\|U_k(a)\|^2\le C\sum_i a_i^2$ hold for all $a$? The answer is that the
rescaled classical constant is already optimal at every level.

\begin{theorem}\label{thm:main}
For all $m\ge2$, $1\le k\le m$, and $a\in\mathbb{R}^m$,
\[
\|U_k(a)\|^2\;\le\;\frac{k(m-k+1)}{m}\,\sum_{i=1}^m a_i^2 .
\]
The constant $C(m,k)=\frac{k(m-k+1)}{m}$ is best possible. Moreover:
\begin{enumerate}
\item[(i)] for $k=1$ and $k=m$ equality holds for every $a$;
\item[(ii)] for $1<k<m$ and $a\ne0$, equality holds if and only if
$|a_1|=\cdots=|a_m|$.
\end{enumerate}
\end{theorem}

The equal-weight case gives $\|U_k(\mathbf 1)\|^2=k(m-k+1)=C(m,k)\cdot m$,
so the theorem asserts that arbitrary coordinate weights, measured in
$\ell^2$, can never beat the rescaled classical constant, and that among
nonzero weights the \emph{only} extremizers in the interior are the
equal-magnitude ones. That a proof beyond spectral methods is
\emph{necessary} is not an accident: for generic $a$ the family
$\{U_k(a)\}$ no longer generates a Bose--Mesner algebra /
$\mathfrak{sl}_2$ action commuting with the Johnson scheme, so the standard
association-scheme diagonalization is unavailable (cf.\ Remark
\ref{rem:spectral}).

The proof occupies \S\ref{sec:proof}: a sign reduction (\S\ref{subsec:sign}),
a coordinate-deletion block decomposition (\S\ref{subsec:block}), a
block-norm comparison lemma (\S\ref{subsec:blocknorm}), and a scalar
one-variable closure inequality whose defect factors as a perfect square
(\S\ref{subsec:closure}). The equality classification is proved in
\S\ref{sec:equality} by tracing equality through every step of the
induction. \S\ref{sec:remarks} records consequences, limitations, and open
questions; Appendix~\ref{app:repro} describes the machine-checked
identities supporting the proof.

\section{Related work}\label{sec:related}

\textbf{Equal-weight spectra.} The spectrum of the unweighted Boolean up
operators, and hence \eqref{eq:classical}, is classical; see
\cite{srinivasan2011} and the association-scheme diagonalization of
\cite{bannai1984}. Our theorem contains \eqref{eq:classical} as the ray
$a=\mathbf 1$ and does not claim it.

\textbf{Inclusion matrices.} Wilson's theorem on inclusion matrices over
fields of characteristic zero determines the \emph{rank} of unweighted
inclusion matrices (see \cite{wilson1984} and the encyclopedic treatment
of Johnson-scheme spectra in \cite{godsilroyle2001}); it is a rank statement, not a norm statement, and it
does not carry weights. Singular-value information for weighted inclusion
matrices does not appear to be tabulated; the present result gives the top
singular value exactly for the Boolean level pairs.

\textbf{Differential posets and $\mathfrak{sl}_2$ calculi.} Stanley's
differential posets \cite{stanley1988} provide weighted up/down calculus
for enumeration identities, not sharp $\ell^2$ operator norms; the
$\mathfrak{sl}_2$-type relations behind \eqref{eq:classical} acquire
non-scalar defect terms for generic $a$ (Remark \ref{rem:spectral}).

\textbf{Analysis of Boolean functions.} Reverse hypercontractivity on the
cube \cite{mos2013}, in the lineage of Bonami--Beckner
\cite{bonami1970,beckner1975}, controls $L^p$--$L^q$ norms of low-degree
polynomials but does not bound a single weighted level-lowering operator in
operator norm; the full-space operator $\sum_i a_i\,\sigma_i^{-}$ on
$\mathbb{R}^{2^{[m]}}$ has norm $\max_k\|U_k(a)\|$, which loses the
level-dependent constant. Linear polarization constants of Hilbert spaces
\cite{matolcsi2005,pappas2004} concern products of linear functionals and
do not imply Theorem \ref{thm:main}.

\section{Preliminaries}\label{sec:prelim}

Throughout, $\|\cdot\|$ denotes the operator norm for the Euclidean norms
on the spaces involved, and we write
$G_k(a):=U_k(a)U_k(a)^{\mathsf T}$ on $\ell^2\binom{[m]}{k-1}$, so that
$\|U_k(a)\|^2=\lambda_{\max}(G_k(a))$.

\subsection{Sign and phase reduction}\label{subsec:sign}

We record the reduction in its exact two-sided form, which also covers
complex weights. Write $a_i=|a_i|e^{i\theta_i}$ ($\theta_i\in\{0,\pi\}$ for
real signs) and define diagonal unitaries $D_r$ on
$\ell^2\binom{[m]}{k-1}$ and $D_c$ on $\ell^2\binom{[m]}{k}$ by
$(D_r e)_S=e^{-i\sum_{j\in S}\theta_j}e_S$ and
$(D_c e)_T=e^{\,i\sum_{j\in T}\theta_j}e_T$. For $T=S\cup\{i\}$ the
$(S,T)$ entry of $D_r\,U_k(|a|)\,D_c$ is
$|a_i|\,e^{-i\sum_{j\in S}\theta_j}e^{\,i\sum_{j\in T}\theta_j}=
|a_i|e^{i\theta_i}=a_i$, so
\begin{equation}\label{eq:phase}
U_k(a)\;=\;D_r\,U_k(|a|)\,D_c ,
\end{equation}
a two-sided unitary equivalence; hence
$\|U_k(a)\|=\|U_k(|a|)\|$ and the equality classification inherits the
reduction. We assume $a\ge0$ from now on.

\subsection{Block decomposition by deleting a coordinate}\label{subsec:block}

Write $a=(a',a_m)$ with $a'\in\mathbb{R}^{m-1}$. Order the rows by
$(k-1)$-sets not containing $m$ first, then those containing $m$, and the
columns likewise. Since the $(S,T)$ entry is nonzero only when
$T=S\cup\{i\}$,
\begin{equation}\label{eq:block}
U_k(a)\;\simeq\;
\begin{pmatrix} U_k(a') & a_m I\\[2pt] 0 & U_{k-1}(a') \end{pmatrix},
\end{equation}
where $I$ is the $\binom{m-1}{k-1}\times\binom{m-1}{k-1}$ identity indexed
by the bijection $S\leftrightarrow S\cup\{m\}$, and $\simeq$ denotes
equality after the stated ordering. Multiplying the block forms, and
writing $A=U_k(a')$, $B=U_{k-1}(a')$ so that
$U_k(a)\simeq\bigl(\begin{smallmatrix}A&a_mI\\0&B\end{smallmatrix}\bigr)$
and
$U_k(a)^{\mathsf T}\simeq\bigl(\begin{smallmatrix}A^{\mathsf T}&0\\a_mI&B^{\mathsf T}\end{smallmatrix}\bigr)$,
\begin{equation}\label{eq:gram}
U_k(a)U_k(a)^{\mathsf T}\;\simeq\;
\begin{pmatrix}
U_k(a')U_k(a')^{\mathsf T}+a_m^2 I & a_m\,U_{k-1}(a')^{\mathsf T}\\
a_m\,U_{k-1}(a') & U_{k-1}(a')U_{k-1}(a')^{\mathsf T}
\end{pmatrix};
\end{equation}
the off-diagonal block is $a_m\,U_{k-1}(a')^{\mathsf T}$ (shape
$\binom{m-1}{k-1}\times\binom{m-1}{k-2}$), so its largest singular value
is \emph{exactly} $a_m\,\|U_{k-1}(a')\|$.

\subsection{A block-norm lemma}\label{subsec:blocknorm}

\begin{lemma}\label{lem:block}
If $M=\bigl(\begin{smallmatrix}P&Q\\ Q^{\mathsf T}&R\end{smallmatrix}\bigr)$
is real symmetric with $P,R$ square symmetric, then
\[
\lambda_{\max}(M)\;\le\;
\lambda_{\max}\begin{pmatrix}\lambda_{\max}(P)&s_1(Q)\\ s_1(Q)&\lambda_{\max}(R)\end{pmatrix},
\]
where $s_1$ is the largest singular value.
\end{lemma}
\begin{proof}
For a unit vector $x=(u,v)$, writing $\alpha=|u|$, $\beta'=|v|$,
\[
x^{\mathsf T}Mx \le \lambda_{\max}(P)\,\alpha^2+\lambda_{\max}(R)\,\beta'^{2}
+2\,s_1(Q)\,\alpha\beta' = y^{\mathsf T}Sy,
\qquad y=(\alpha,\beta'),\ |y|=1,
\]
by Cauchy--Schwarz applied to $u^{\mathsf T}Qv$. Maximizing over unit $y$
gives $\lambda_{\max}(S)$.
\end{proof}

\section{Proof of Theorem \ref{thm:main}}\label{sec:proof}

\subsection{The induction}\label{subsec:induction}

We prove the theorem, simultaneously in $k$, by induction on $m$, carrying
the equality clause. \emph{Induction hypothesis $\mathrm{IH}(m')$} for
$m'<m$: the theorem holds for all $1\le k'\le m'$ and all weights in
$\mathbb{R}^{m'}$. The base $m=2$ is the endpoint case below. The endpoints
for general $m$ are exact: $U_1(a)$ is the row vector $(a_1,\dots,a_m)$ and
$U_m(a)$ a column vector, so
$\|U_1(a)\|^2=\|U_m(a)\|^2=|a|^2=C(m,1)|a|^2=C(m,m)|a|^2$, and every $a$
is an extremizer, proving (i).

Let $1<k<m$ and $a\ne0$. If $a'=0$ then
$U_k(a)=a_m$ times a partial isometry given by the bijection, so
$\|U_k(a)\|^2=a_m^2<C(m,k)\,a_m^2$ because
$k(m-k+1)-m=(k-1)(m-k)>0$; this also shows no extremizer has exactly one
nonzero coordinate. Assume then $s:=|a'|^2>0$ and set
$t=a_m^2/s$, $\alpha^2=\|U_k(a')\|^2$, $\beta^2=\|U_{k-1}(a')\|^2$.
Applying Lemma \ref{lem:block} to \eqref{eq:gram}, with the exact identity
$s_1\bigl(a_mU_{k-1}(a')^{\mathsf T}\bigr)=a_m\beta$,
\[
\|U_k(a)\|^2\;\le\;\lambda_{\max}
\begin{pmatrix}\alpha^2+a_m^2& a_m\beta\\ a_m\beta&\beta^2\end{pmatrix}.
\]
By $\mathrm{IH}(m-1)$,
$\alpha^2\le A\,s$ and $\beta^2\le B\,s$ with
$A=C(m-1,k)=\frac{k(m-k)}{m-1}$ and $B=C(m-1,k-1)=\frac{(k-1)(m-k+1)}{m-1}$;
both children lie in range because $1\le k-1$ and $k\le m-1$.
Since $\lambda_{\max}\bigl(\begin{smallmatrix}p&o\\ o&q\end{smallmatrix}\bigr)
=\tfrac12\bigl(p+q+\sqrt{(p-q)^2+4o^2}\bigr)$ is nondecreasing in each of
$p$, $q$ and in $|o|$ (the partial derivatives have the form
$\tfrac12[1\pm(p-q)/\sqrt{(p-q)^2+4o^2}]\ge0$, and the bound in $o$ is
even),
\begin{equation}\label{eq:corner}
\|U_k(a)\|^2\;\le\; s\cdot
\lambda_{\max}\begin{pmatrix}A+t&\sqrt{tB}\\ \sqrt{tB}&B\end{pmatrix}.
\end{equation}

\subsection{Scalar closure}\label{subsec:closure}

\begin{lemma}\label{lem:closure}
For $m\ge3$, $1\le k\le m$, $t\ge0$, with $A,B$ as above and
$C=C(m,k)=\frac{k(m-k+1)}{m}$,
\[
\lambda_{\max}\begin{pmatrix}A+t&\sqrt{tB}\\ \sqrt{tB}&B\end{pmatrix}
\;\le\; C(1+t),
\]
and for interior $k$ equality holds if and only if $t=\frac1{m-1}$.
\end{lemma}
\begin{proof}
The largest eigenvalue equals $\tfrac12\bigl(A+B+t+\sqrt{(A-B+t)^2+4Bt}\bigr)$,
so the claim is $\sqrt{(A-B+t)^2+4Bt}\le X(t)$ with
$X(t)=2C(1+t)-A-B-t$. Now $X$ is affine with slope
$2C-1=\frac{2k(m-k+1)-m}{m}>0$ and intercept
$X(0)=2C-A-B=\frac{k(k-1)+(m-k)(m-k+1)}{m(m-1)}>0$; hence $X(t)>0$ for
$t\ge0$ and squaring is legitimate. The squared defect factors exactly:
\[
X(t)^2-\bigl[(A-B+t)^2+4Bt\bigr]
=\frac{4k(k-1)(m-k)(m-k+1)}{m^2(m-1)^2}\,\bigl((m-1)t-1\bigr)^2\ \ge\ 0,
\]
a polynomial identity in $m,k,t$ verified symbolically (see
Appendix~\ref{app:repro}). For interior $k$
the prefactor is positive, so equality holds iff $t=\frac1{m-1}$.
\end{proof}

Combining \eqref{eq:corner} with Lemma \ref{lem:closure},
$\|U_k(a)\|^2\le s\,C(1+t)=C(m,k)(s+a_m^2)=C(m,k)|a|^2$, completing the
induction. Sharpness: $G_k(\mathbf 1)\mathbf 1=k(m-k+1)\mathbf 1$, so
equality is attained at $a=\mathbf 1$ in every level and no smaller
constant works at any level. \hfill$\square$

\begin{corollary}\label{cor:complex}
Theorem \ref{thm:main} remains true for $a\in\mathbb{C}^m$ with
$\sum_i|a_i|^2$ on the right, and the equality classification is unchanged
with $|a_i|$ in place of $a_i$.
\end{corollary}
\begin{proof}
By \eqref{eq:phase}, $\|U_k(a)\|=\|U_k(|a|)\|$ entrywise in magnitudes, and
$\sum_i a_i^2$ is replaced by $\sum_i|a_i|^2$ throughout; equality holds
for $a$ exactly when it holds for $|a|$.
\end{proof}

\section{Equality cases}\label{sec:equality}

Assume $1<k<m$ and equality in Theorem \ref{thm:main} with $a\ne0$.
Tracing the chain in \S\ref{sec:proof}: Lemma \ref{lem:closure} must be an
equality, forcing $t=\frac1{m-1}$; since then $t>0$ and $B>0$ the corner
off-diagonal $\sqrt{tB}\,s$ is nonzero, and the monotonicity step
\eqref{eq:corner} is strict in each diagonal entry when the off-diagonal
is nonzero (the partial derivatives of $\lambda_{\max}$ are
$\frac12\bigl[1\pm\frac{p-q}{\sqrt{(p-q)^2+4o^2}}\bigr]>0$ for $o>0$,
because $\sqrt{(p-q)^2+4o^2}>|p-q|$), forcing $\alpha^2=As$ and
$\beta^2=Bs$. Thus equality holds
at the subproblem $(m-1,k)$ on $a'$ (and at $(m-1,k-1)$). For $m\ge4$ at
least one of these subproblems is interior --- if $2\le k\le m-2$ then
$(m-1,k)$ is; if $k=2$ (resp.\ $k=m-1$) then $(m-1,2)$ (resp.\
$(m-1,m-2)$) is interior for $m\ge4$ --- and the strengthened induction
hypothesis yields $|a_1|=\cdots=|a_{m-1}|$; then $t=\frac1{m-1}$ gives
$a_m^2=s/(m-1)$, the same common value. The remaining base $m=3,k=2$ is
computed exactly: with $x_i=a_i^2$ and $s=x_1+x_2+x_3$, the matrix
$\frac43 s\,I-G_2(a)$ has principal minors
$\frac{x_1+x_2+4x_3}{3}$, $\frac{s(x_1+4x_2+4x_3)}{9}$; these are the $(3,3)$
entry and the $2\times2$ principal minor on indices $\{2,3\}$ in the basis
ordering $e_{\{1\}},e_{\{2\}},e_{\{3\}}$, and the remaining principal minors
are obtained by permuting the coordinate labels $1,2,3$,
and determinant $\frac4{27}\bigl[(x_1+x_2+x_3)^3-27x_1x_2x_3\bigr]\ge0$ by
AM--GM, with equality iff $x_1=x_2=x_3$. Hence
$\lambda_{\max}(G_2(a))=\frac43 s$ iff $a_1^2=a_2^2=a_3^2$. The converse
direction (equal magnitudes give equality) follows from the sign reduction
\eqref{eq:phase} and $G_k(\mathbf 1)\mathbf 1=k(m-k+1)\mathbf 1$.
This completes the proof of Theorem \ref{thm:main}(ii). \hfill$\square$

\begin{remark}[complement duality]\label{rem:duality}
Let $C_j:\ell^2\binom{[m]}{j}\to\ell^2\binom{[m]}{m-j}$, $e_S\mapsto
e_{S^{\mathsf c}}$. Since $(S\cup\{i\})^{\mathsf c}=S^{\mathsf c}\setminus\{i\}$,
entrywise comparison gives $C_{k-1}\,U_k(a)\,C_k=U_{m-k+1}(a)^{\mathsf T}$,
hence $\|U_k(a)\|=\|U_{m-k+1}(a)\|$ and $C(m,k)=C(m,m-k+1)$; the induction
above therefore needs only the range $k\le\frac{m+1}{2}$.
\end{remark}

\begin{remark}[spectral methods fail for generic weights]\label{rem:spectral}
For $a=\mathbf 1$ the operators $\{U_k,U_k^{\mathsf T}\}$ act between the
homogeneous components of the Boolean lattice with
$[D,U]=U$-type relations after weighting by rank, and the spectra follow
from the Johnson scheme \cite{srinivasan2011}. For generic $a$ the
commutation relations acquire non-scalar defect terms and no association
scheme survives; the deletion induction appears to be the only route to a
sharp general result.
\end{remark}

\section{Concluding remarks}\label{sec:remarks}

\textbf{Consequences.} Theorem \ref{thm:main} identifies the exact worst
case of a weighted level-lowering operator on the Boolean lattice: the
worst weights are the flat ones, up to signs and phases
(Corollary~\ref{cor:complex}), and the extremal rigidity is exact rather
than asymptotic. By complement duality
(Remark~\ref{rem:duality}) the same bound and classification hold for the
weighted down-operators $U_k(a)^{\mathsf T}$. The scalar closure lemma is
an independent one-variable inequality that may be reusable wherever a
$2\times2$ corner comparison of this shape arises.

\textbf{Limitations.} The method is specific to the Boolean lattice level
structure: the deletion recursion produces exactly two children with the
same ambient parameters, and the perfect-square factorization depends on
the constants $C(m,k)$. No comparison with the weighted spectra of
$G_k(a)$ beyond the top eigenvalue is obtained, and nothing in this paper
addresses the infinite-dimensional operator-theoretic problems
(e.g.\ Nyman--Beurling-type formulations) that motivated the search; any
connection there is conditional and unproved.

\textbf{Open questions.} (1) For which graded posets with rank numbers
analogous to $\binom{mj}{j}$ does the deletion induction close with a
perfect-square defect? (2) Determine $\lambda_{\max}(G_k(a))$ exactly as a
function of the sorted magnitudes of $a$, not only its maximum over the
sphere $\{|a|^2=1\}$. (3) Analogous sharp bounds for the two-step
operators $U_k(a)U_{k+1}(a)$, where the corner comparison no longer
closes at one variable.

\begin{thebibliography}{9}
\bibitem{srinivasan2011}
E.~Srinivasan, \emph{Counting spanning trees of almost regular graphs and
block diagonalization}, arXiv:1104.1481; and references therein for the
spectrum of the Boolean up operators.
\bibitem{bannai1984}
E.~Bannai and T.~Ito, \emph{Algebraic Combinatorics. I. Association
Schemes}, Benjamin/Cummings Publishing Co., Menlo Park, CA, 1984.
\bibitem{mos2013}
E.~Mossel, K.~Oleszkiewicz and A.~Sen, \emph{On reverse hypercontractivity},
Geom. Funct. Anal.\ \textbf{23} (2013), 1062--1097. arXiv:1108.1210.
\bibitem{stanley1988}
R.~P.~Stanley, \emph{Differential posets}, J. Amer. Math. Soc.\ \textbf{1}
(1988), 919--961.
\bibitem{matolcsi2005}
M.~Matolcsi, \emph{The linear polarization constant of $\mathbb{R}^n$},
Acta Math. Hungar.\ \textbf{108} (2005), 129--136.
\bibitem{pappas2004}
A.~Pappas and Sz.~Gy.~R\'ev\'esz, \emph{Linear polarization constants of
Hilbert spaces}, J. Math. Anal. Appl.\ \textbf{300} (2004), 129--146.
\bibitem{bonami1970}
A.~Bonami, \emph{Ensembles $\Lambda_p$ dans le produit d'espaces de Banach
et application}, C.\ R.\ Acad.\ Sci.\ Paris S\'er.\ A--B\ \textbf{270}
(1970), 1335--1338.
\bibitem{beckner1975}
W.~Beckner, \emph{Inequalities in Fourier analysis}, Ann.\ of Math.\ \textbf{102} (1975), 159--182.
\bibitem{wilson1984}
R.~M.~Wilson, \emph{The exact bound in the Erd\H{o}s--Ko--Rado theorem},
Combinatorica \textbf{4} (1984), 247--257.
\bibitem{godsilroyle2001}
C.~D.~Godsil and G.~F.~Royle, \emph{Algebraic Graph Theory}, Graduate
Texts in Mathematics \textbf{207}, Springer, New York, 2001.
\end{thebibliography}

\appendix
\section{Reproducibility of the checked identities}\label{app:repro}

The proof above contains no step that requires computer assistance; the
machine checks below are regression evidence, not part of the proof.

\textbf{Closure identity.} With $A,B,C$ as in Lemma \ref{lem:closure}, the
defect factorization
\[
X(t)^2-\bigl[(A-B+t)^2+4Bt\bigr]
=\frac{4k(k-1)(m-k)(m-k+1)}{m^2(m-1)^2}\bigl((m-1)t-1\bigr)^2
\]
was verified by expansion and simplification in \texttt{sympy} as a
rational-function identity in \emph{general} symbols $(m,k,t)$, not by
sampling.

\textbf{Independent reproduction.} A verification script built from the
definition of $U_k(a)$ alone (no reuse of the proof) checks the block
identity \eqref{eq:block}, the Gram identity \eqref{eq:gram}, the duality
identity of Remark \ref{rem:duality}, and the phase identity
\eqref{eq:phase} entrywise for small parameters, and evaluates the
inequality at $30$ exact rational instances $(m,k,a)$ (worst observed
ratio $0.979016$, zero violations), using certified real-root isolation
for the spectral maxima. Exact checks of $\lambda_{\max}(U_k(\mathbf 1)
U_k(\mathbf 1)^{\mathsf T})=k(m-k+1)$ were run for $2\le m\le 8$, all $k$.

\textbf{Artifacts.} The scripts producing the checks described above,
with pinned dependency versions and reproduction instructions, are
provided in the \texttt{verification/} directory of the accompanying
repository (\texttt{weighted-johnson-inequality}); the symbolic script
re-derives the closure identity in general symbols and the exact script
rebuilds $U_k(a)$ from its definition alone. No mathematical step of this
manuscript depends on any of these machine checks.

\end{document}
